\documentclass[11pt]{article}
\usepackage[margin=1in]{geometry}
\usepackage{enumitem}
% =============================================================================
% Preambles/header.tex — shared LaTeX/Beamer preamble for this project
%
% Use from a Beamer lecture by prepending:
%     \input{header}
% and compiling with TEXINPUTS=../Preambles:$TEXINPUTS (the /compile-latex
% skill does this automatically).
%
% PALETTE CONTRACT. Color names defined here MUST match the names in
% Quarto/theme-template.scss so Beamer and Quarto renderings use the same
% palette. The scripts/check-palette-sync.sh script enforces this.
%
% To customize: edit the HEX values below AND the matching $variable in
% Quarto/theme-template.scss, then run scripts/check-palette-sync.sh.
% =============================================================================

% -----------------------------------------------------------------------------
% Engine detection + encoding
%
% Our /compile-latex skill uses XeLaTeX, where inputenc is unnecessary and
% can produce encoding warnings. Only load inputenc under pdfTeX.
% -----------------------------------------------------------------------------
\usepackage{iftex}
\ifPDFTeX
  \usepackage[utf8]{inputenc}
\fi

\usepackage{amsmath, amssymb, amsthm}
\usepackage{booktabs}
\usepackage{graphicx}

% -----------------------------------------------------------------------------
% PALETTE — mirrors Quarto/theme-template.scss variable names.
% Load xcolor and define colors BEFORE anything that references them
% (notably hyperref, hypersetup, and the Beamer color assignments).
% -----------------------------------------------------------------------------
\usepackage{xcolor}

\definecolor{primary-blue}{HTML}{012169}      % headings, accents
\definecolor{primary-gold}{HTML}{B9975B}      % emphasis, borders
\definecolor{highlight-yellow}{HTML}{F2A900}  % markers, alerts
\definecolor{light-bg}{HTML}{E8EDF5}          % title / section backgrounds
\definecolor{jet}{HTML}{1A1A1A}               % body text

% Semantic colors (used in diagrams and annotations)
\definecolor{positive}{HTML}{15803D}          % good / observed / identified
\definecolor{negative}{HTML}{B91C1C}          % bad / problematic / bias
\definecolor{neutral}{HTML}{525252}           % reference / context

% Highlight accents (match .hi-* classes in the SCSS)
\definecolor{hi-slate}{HTML}{314F4F}
\definecolor{hi-green}{HTML}{2E8B57}
\definecolor{hi-red}{HTML}{C41E3A}

% Semantic aliases so skill templates and snippets can write
% \textcolor{accent}{...} without hard-coding "primary-blue".
\colorlet{accent}{primary-blue}
\colorlet{accent2}{primary-gold}

% -----------------------------------------------------------------------------
% Hyperref — loaded AFTER xcolor + palette so link colors resolve.
% -----------------------------------------------------------------------------
\usepackage{hyperref}
\hypersetup{colorlinks=true, linkcolor=primary-blue, urlcolor=primary-blue,
            citecolor=primary-blue, pdfborder={0 0 0}}

% -----------------------------------------------------------------------------
% Course code — set once per deck (after \input{header}, before \begin{document})
% via \coursecode{CS401}. Shown in the footer on every slide so a deck's course
% is identifiable without opening the title page. Empty by default.
% -----------------------------------------------------------------------------
\makeatletter
\newcommand{\coursecode}[1]{\gdef\@coursecodetext{#1}}
\gdef\@coursecodetext{}
\makeatother

% -----------------------------------------------------------------------------
% Beamer theme assignments (only apply under Beamer).
%
% We detect Beamer via \@ifclassloaded{beamer}, which is the documented,
% non-internal way. This must be wrapped in \makeatletter/\makeatother
% because \@ifclassloaded contains an @-catcode character.
% -----------------------------------------------------------------------------
\makeatletter
\@ifclassloaded{beamer}{%
  \usefonttheme{professionalfonts}
  \setbeamercolor{structure}{fg=primary-blue}
  \setbeamercolor{frametitle}{fg=primary-blue}
  \setbeamercolor{title}{fg=primary-blue}
  \setbeamercolor{subtitle}{fg=primary-gold}
  \setbeamercolor{author}{fg=primary-blue}
  \setbeamercolor{institute}{fg=neutral}
  \setbeamercolor{date}{fg=primary-gold}
  \setbeamercolor{itemize item}{fg=primary-blue}
  \setbeamercolor{itemize subitem}{fg=primary-gold}
  \setbeamercolor{alerted text}{fg=primary-gold}
  \setbeamercolor{block title}{fg=white, bg=primary-blue}
  \setbeamercolor{block body}{bg=light-bg}
  % Semantic block variants: block=definition/concept (blue), exampleblock=
  % worked example (green/positive), alertblock=key takeaway or caution (gold).
  % Keep to <=2 colored boxes per slide across ALL three variants (INV-7).
  \setbeamercolor{block title example}{fg=white, bg=positive}
  \setbeamercolor{block body example}{bg=light-bg}
  \setbeamercolor{block title alerted}{fg=white, bg=primary-gold}
  \setbeamercolor{block body alerted}{bg=light-bg}

  % Minimal footer: course code (left) + page X / N (right), no navigation clutter
  \setbeamertemplate{navigation symbols}{}
  \setbeamertemplate{footline}{%
    \color{neutral}\footnotesize\hspace{0.7em}\@coursecodetext\hfill
    \insertframenumber/\inserttotalframenumber\hspace{0.7em}\vspace{0.3em}}
}{}
\makeatother

% -----------------------------------------------------------------------------
% TikZ setup — libraries the snippet gallery uses
% -----------------------------------------------------------------------------
\usepackage{tikz}
\usetikzlibrary{arrows.meta, positioning, calc, decorations.pathreplacing,
                fit, shapes.geometric, backgrounds}

% Shared TikZ styles — snippets and hand-written diagrams can pick these up
% via `\begin{tikzpicture}[every node/.style={...}]` or by naming specific
% styles (e.g., `\node[dag-node] (X) at (0,0) {X};`).
\tikzset{
  % Node styles for causal DAGs and similar
  dag-node/.style = {
    draw=primary-blue, fill=light-bg, rounded corners=2pt,
    minimum width=1.2cm, minimum height=0.9cm, align=center, font=\small
  },
  decision-node/.style = {
    draw=primary-gold, fill=white, diamond, aspect=1.8,
    minimum width=1.8cm, minimum height=1.2cm, align=center, font=\small,
    inner sep=1pt
  },
  % Edge styles — observed vs counterfactual
  observed-edge/.style = {-{Latex[length=2.5mm]}, thick, draw=jet},
  counterfactual-edge/.style = {-{Latex[length=2.5mm]}, thick, draw=neutral, dashed},
  confound-edge/.style = {-{Latex[length=2.5mm]}, thick, draw=negative, dashed},
  % Data-point styles for plots
  observed-dot/.style = {fill=primary-blue, draw=primary-blue, circle, inner sep=1.2pt},
  counterfactual-dot/.style = {fill=white, draw=neutral, circle, inner sep=1.2pt}
}

% -----------------------------------------------------------------------------
% Convenience macros
% -----------------------------------------------------------------------------
\newcommand{\muted}[1]{\textcolor{neutral}{#1}}
\newcommand{\key}[1]{\textcolor{primary-gold}{\textbf{#1}}}
\newcommand{\good}[1]{\textcolor{positive}{#1}}
\newcommand{\bad}[1]{\textcolor{negative}{#1}}

% Standout frame macro: full-bleed dark background for section transitions.
% Beamer-only (uses \begin{frame}/\setbeamercolor/\usebeamerfont) -- do not
% call from an article-class file (e.g. Notes/<CODE>/*-notes.tex). \muted,
% \key, \good, \bad, and \coursecode above are plain xcolor/gdef and are
% safe to use from either class; verified when Notes/article-class support
% was added.
% Expand: \transitionslide{Your transition title}
\newcommand{\transitionslide}[1]{%
  \begingroup
  \setbeamercolor{background canvas}{bg=primary-blue}
  \begin{frame}[plain]
    \centering
    \vfill
    {\usebeamerfont{title}\color{white}\Large #1\par}
    \vfill
  \end{frame}
  \endgroup
}

% Section divider frame (Beamer-only), an alternative to \transitionslide with a
% darker two-line style: near-black (jet) background, a small white label line
% above a larger gold title, and a thin gold rule along the bottom edge.
% Expand: \sectiondivider{Section 2}{Pointers}
\newcommand{\sectiondivider}[2]{%
  \begingroup
  \setbeamercolor{background canvas}{bg=jet}
  \begin{frame}[plain]
    \vspace*{3.0cm}
    {\color{white}\ttfamily\large #1\par}
    \vspace{0.5cm}
    {\color{primary-gold}\LARGE #2\par}
    \begin{tikzpicture}[remember picture, overlay]
      \draw[primary-gold, line width=2.5pt]
        ($(current page.south west)+(0,0.1)$) -- ($(current page.south east)+(0,0.1)$);
    \end{tikzpicture}
  \end{frame}
  \endgroup
}

% -----------------------------------------------------------------------------
% Channel watermark — "Concepts That Click" (YouTube).
%
% Article-class files (CompetitiveExam Q/A sets, Notes, Assignments) get a
% light diagonal draft watermark on every page plus a centered footer naming
% the channel. Beamer decks get a subtle TikZ background watermark instead
% (the footline already carries the course code). Centralized here so every
% MCQ PDF is branded without per-file edits.
% -----------------------------------------------------------------------------
\makeatletter
\@ifclassloaded{article}{%
  \usepackage{draftwatermark}
  \SetWatermarkText{Concepts That Click}
  \SetWatermarkScale{1}
  \SetWatermarkFontSize{2cm}
  \SetWatermarkLightness{0.86}
  \SetWatermarkAngle{45}
  \usepackage{fancyhdr}
  \pagestyle{fancy}
  \fancyhf{}
  \fancyfoot[C]{\footnotesize\textcolor{neutral}{Concepts That Click~|~YouTube: \href{https://www.youtube.com/@concepts.thatclick}{@concepts.thatclick}}}
  \renewcommand{\headrulewidth}{0pt}
  \renewcommand{\footrulewidth}{0pt}
  \fancypagestyle{plain}{%
    \fancyhf{}%
    \fancyfoot[C]{\footnotesize\textcolor{neutral}{Concepts That Click~|~YouTube: \href{https://www.youtube.com/@concepts.thatclick}{@concepts.thatclick}}}%
    \renewcommand{\headrulewidth}{0pt}%
    \renewcommand{\footrulewidth}{0pt}%
  }
}{}
\@ifclassloaded{beamer}{%
  \setbeamertemplate{background}{%
    \begin{tikzpicture}[remember picture, overlay]
      \node[rotate=45, opacity=0.055, align=center]
        at (current page.center)
        {\Huge\textcolor{neutral}{Concepts That Click}};
    \end{tikzpicture}%
  }
}{}
\makeatother 
\coursecode{JKSSB}
\renewcommand{\thesection}{30.\arabic{section}}
\newtheorem{example}{Example}[section]
\newenvironment{definitionbox}[1]{%
  \par\noindent\textbf{Definition (#1).}\ \itshape
}{\par\vspace{0.3em}}
\title{Access Queries, Forms and Reports --- Detailed Lecture Notes}
\author{JKSSB Computer Awareness}
\date{\today}

\begin{document}
\maketitle

\section{The Access objects and the object sequence}
An Access database is built from a small set of \textbf{objects}. The
\textbf{table} stores the actual data. The \textbf{query} asks a question
and retrieves selected data. The \textbf{form} provides a screen for
entering and viewing data. The \textbf{report} presents data in a formatted,
printable layout. These four objects are used together in a natural order:
data is stored in a table, selected by a query, entered through a form, and
printed by a report. That order --- \textbf{Table} $\rightarrow$
\textbf{Query} $\rightarrow$ \textbf{Form} $\rightarrow$ \textbf{Report}
--- is the object sequence this lesson is built around.

\begin{definitionbox}{Database object}
A named component of an Access database that performs one job. The main
objects are the table, query, form and report.
\end{definitionbox}

\begin{center}
\begin{tabular}{p{0.16\textwidth}p{0.68\textwidth}}
\toprule
\textbf{Object} & \textbf{Job} \\
\midrule
Table & Stores the actual data in rows and columns. \\
Query & Asks a question and retrieves selected data. \\
Form & Provides an on-screen interface to enter and view data. \\
Report & Presents data in a formatted, printable layout. \\
\bottomrule
\end{tabular}
\end{center}

The rest of this lesson uses these terms in one consistent sense. A
\textbf{table} is the store. A \textbf{query} is the question. A
\textbf{form} is the on-screen entry screen. A \textbf{report} is the printed
output.

\section{The table is the source of data}
A table stores data in \textbf{fields} (columns) and \textbf{records}
(rows). Every other object draws its data from a table or from a query. A
form or report built on a table shows that table's data; built on a query,
it shows only the rows the query selected. No data is stored in a query, a
form or a report --- those objects only \emph{present} what the table holds.
This is why the table is called the foundation of the database.

\section{What a query is}
A \textbf{query} is a question asked of the database. It retrieves data from
one or more tables, usually by a condition. A query does not normally store
new data of its own; it shows a \emph{result set} drawn from the tables each
time it is opened. Queries can also be used to \emph{change} data; those are
called \textbf{action queries}, and they are discussed below.

\begin{definitionbox}{Query}
An Access object that selects, filters, or acts on records according to the
rules the user sets, drawing its data from one or more tables.
\end{definitionbox}

\section{The Select query: retrieving on a condition}
A \textbf{Select query} retrieves the records that satisfy a
\textbf{condition} (also called a \emph{criterion}). It is the most common
query type in Access. The condition decides \emph{which rows appear}, and
the chosen fields decide \emph{which columns appear}. For example, a Select
query on a Students table could show only those students whose
\texttt{Marks} are greater than 60.

\begin{example}
An item asks: ``What is a query that retrieves data on the basis of a
condition called?'' The query is only \emph{retrieving} data --- it is not
changing it --- so it is a \textbf{Select query}. This is the concept tested
by PYQ-39 (Website Operator 2026, query retrieving data on a condition).
\end{example}

\section{Building a query: Design view and criteria}
Queries are usually built in \textbf{Query Design} view, or typed directly in
\textbf{SQL} view. The design grid is called \textbf{QBE} ---
\textbf{Query By Example}. Each column of the grid holds one \textbf{field},
and the \textbf{Criteria} row holds the condition for that field. A criterion
of \texttt{>60} on the \texttt{Marks} field, for instance, keeps only the
rows in which marks are above 60.

\begin{definitionbox}{QBE (Query By Example)}
The grid-based design view in Access in which a query is built by choosing
fields and typing criteria, without writing the SQL by hand.
\end{definitionbox}

\begin{example}
For a Students table, place \texttt{Marks} in the field row and \texttt{>60}
in the Criteria row. When the query runs, Access returns only the students
who scored more than 60.
\end{example}

\section{Action queries: Update, Append and Delete}
An \textbf{action query} changes the stored data. The three common types are:
\begin{itemize}
  \item An \textbf{Update query} changes the existing values in records.
  \item An \textbf{Append query} adds records from one table to another.
  \item A \textbf{Delete query} removes the records that match a condition.
\end{itemize}
Because these queries alter the table itself, they are treated with care in
Access. The distinction between retrieving data and changing data is the key
one: a Select query only reads, while Update, Append and Delete queries
write.

\begin{center}
\begin{tabular}{p{0.18\textwidth}p{0.68\textwidth}}
\toprule
\textbf{Query} & \textbf{What it does} \\
\midrule
Select & Retrieves records that satisfy a condition; does not change data. \\
Update & Changes existing values in records. \\
Append & Adds records from one table to another. \\
Delete & Removes records that match a condition. \\
\bottomrule
\end{tabular}
\end{center}

\section{Select versus action queries}
The nearest confusion in this topic is between a Select query and the action
queries. A \textbf{Select query} \textbf{retrieves} data on a condition and
leaves the tables unchanged. An \textbf{action query} (Update, Append,
Delete) \textbf{modifies} the stored data. If a question asks which query
\emph{changes} records, the answer is an action query; if it asks which query
\emph{shows} records that meet a condition, the answer is a Select query.
This is the distinction recorded as TRAP-31.

\begin{example}
An item lists Select, Update, Append and Delete queries and asks which one
does \textbf{not} modify the data. Select is the only one that leaves the
tables unchanged, so Select is the answer.
\end{example}

\section{Lookup fields}
A \textbf{lookup} field lets the user \emph{pick a value from a list} instead
of typing it. The list usually comes from a table or a query --- for example,
a list of department names from a Departments table. A lookup is often used
for a \textbf{foreign key}, so that the value entered matches an existing
value in the related table. The benefit is fewer typing errors and more
consistent data. A lookup field only \emph{chooses} a value from a list; it
does not create a new table or a new query.

\begin{definitionbox}{Lookup field}
A field that displays a list of values, drawn from a table or query, so the
user can select one instead of typing it.
\end{definitionbox}

\section{Basic SQL awareness}
\textbf{SQL} stands for \textbf{Structured Query Language}. It is the
standard language for working with relational databases, and Access runs SQL
behind the scenes whenever a query is run. A retrieval statement has the
general shape \texttt{SELECT fields FROM table WHERE condition;}. The
\texttt{SELECT} keyword names the fields to show, \texttt{FROM} names the
table, and \texttt{WHERE} gives the condition. To retrieve only the students
with marks above 60, one would write
\texttt{SELECT Name FROM Students WHERE Marks > 60;}.

\begin{definitionbox}{SQL}
Structured Query Language --- the standard language used to define and
retrieve data in a relational database.
\end{definitionbox}

\section{Forms: entering and viewing data}
A \textbf{form} is an object that provides an on-screen interface for
entering, viewing and editing data. A form is usually based on a table or a
query. It shows one record at a time (or a few), with clearly labelled
fields, and makes data entry easier and safer than typing directly into a
table grid. Where the table is the raw store, the form is the friendly
screen the user actually works with.

\section{Reports: printable output}
A \textbf{report} presents data in a formatted, printable layout. It can
group records and show totals and summaries. A report is \textbf{read-only}:
it is meant for printing and sharing, not for editing data. Like a form, a
report can be based on a table or a query. The report is the output object of
the database.

\begin{example}
A school can keep all student marks in a table, use a Select query to pick
the students who passed, enter new students through a form, and print the
pass list as a report. The same data moves, in order, through the
Table~$\rightarrow$~Query~$\rightarrow$~Form~$\rightarrow$~Report sequence.
\end{example}

\section{Forms versus reports}
The nearest confusion between the two presentation objects is worth stating
plainly. A \textbf{form} is used \emph{on screen} to enter and view data, and
it allows editing. A \textbf{report} is \emph{read-only} and is used to
\emph{print} data and present summaries. A form answers the question ``how do
I put data in and look at it?''; a report answers ``how do I print it?''

\begin{center}
\begin{tabular}{p{0.18\textwidth}p{0.36\textwidth}p{0.34\textwidth}}
\toprule
\textbf{Point} & \textbf{Form} & \textbf{Report} \\
\midrule
Purpose & Enter and view data on screen & Print and present data \\
Data change & Allows editing & Read-only \\
Typical use & Data-entry screen & Printed summary \\
\bottomrule
\end{tabular}
\end{center}

\section{Exam retrieval and common confusions}
Six distinctions clear up most errors on this topic.
\begin{enumerate}
  \item A \textbf{Select} query \textbf{retrieves} data on a condition;
    \textbf{Update}, \textbf{Append} and \textbf{Delete} queries
    \textbf{modify} data (TRAP-31).
  \item Data is stored in a \textbf{table}, \textbf{not} in a query, form or
    report.
  \item A \textbf{form} is for on-screen entry and viewing; a \textbf{report}
    is a printable, read-only output.
  \item The object sequence is \textbf{Table} $\rightarrow$ \textbf{Query}
    $\rightarrow$ \textbf{Form} $\rightarrow$ \textbf{Report}.
  \item \textbf{SQL} stands for \textbf{Structured Query Language};
    \texttt{SELECT ... FROM ... WHERE} retrieves data on a condition.
  \item A \textbf{lookup} field picks a value from a list supplied by a table
    or query.
\end{enumerate}

\begin{example}
An item states, ``A report is used to enter new records into the
database.'' This is wrong. Entering records is the job of a \textbf{form}; a
report is \textbf{read-only} and is used for printing. The wrong option has
swapped the form and the report.
\end{example}

\subsection*{Exercises}
\begin{enumerate}[label=\arabic*.]
  \item Name the four main objects of an Access database and state the job of
    each.
  \item Give the correct object sequence for storing, selecting, entering and
    printing data.
  \item What is a Select query, and how does it differ from an Update query?
  \item Name the three common action queries and say what each one does.
  \item What is a lookup field, and why is it useful?
  \item Expand SQL and write the general form of a retrieval statement.
  \item Distinguish a form from a report.
\end{enumerate}

\subsection*{Solutions}
\begingroup\small
\begin{enumerate}[label=\arabic*.]
  \item \textbf{Table} stores the data; \textbf{Query} asks a question and
    retrieves selected data; \textbf{Form} provides an on-screen interface to
    enter and view data; \textbf{Report} presents data in a formatted,
    printable layout.
  \item \textbf{Table} $\rightarrow$ \textbf{Query} $\rightarrow$
    \textbf{Form} $\rightarrow$ \textbf{Report}. The table stores the data,
    the query selects it, the form enters it, and the report prints it.
  \item A Select query \textbf{retrieves} records that satisfy a condition
    and leaves the tables unchanged. An Update query \textbf{changes} the
    existing values in records, so it modifies the stored data.
  \item \textbf{Update} changes existing values in records;
    \textbf{Append} adds records from one table to another;
    \textbf{Delete} removes records that match a condition.
  \item A lookup field displays a list of values, drawn from a table or
    query, so the user selects one instead of typing it. It reduces typing
    errors and keeps values consistent, and it is often used for a foreign
    key.
  \item SQL is \textbf{Structured Query Language}. A retrieval statement has
    the form \texttt{SELECT fields FROM table WHERE condition;}, where
    \texttt{SELECT} names the fields, \texttt{FROM} names the table, and
    \texttt{WHERE} gives the condition.
  \item A \textbf{form} is an on-screen interface used to enter and view
    data, and it allows editing. A \textbf{report} is a read-only, formatted
    output used for printing and summarising data.
\end{enumerate}
\endgroup
\end{document}
